\documentclass[10pt]{article}
\usepackage[T1]{fontenc}
\usepackage[utf8]{inputenc}
\usepackage[french]{babel}
\usepackage{amsmath,amssymb,mathrsfs,mathpazo}
\usepackage{geometry}
\geometry{a4paper,margin=2cm}
\usepackage{graphicx,array,tabularx,tabularray,colortbl,multirow,multicol,booktabs,enumitem,pifont,ulem,fancybox,diagbox,cancel}
\usepackage{pstricks,pstricks-add,pst-plot,pst-tree,pst-node,pst-text,pst-func,pst-math,pst-bspline,pst-3dplot,pst-eucl,pst-solides3d}
\let\pscalculate\undefined
\usepackage{pst-bezier}
\usepackage{tikz,pgfplots,tkz-tab}
\usetikzlibrary{arrows,arrows.meta,patterns,calc,positioning,trees,decorations.markings,shapes.geometric,shapes.misc,shapes.arrows,intersections,decorations.pathmorphing,decorations.pathreplacing,backgrounds,shadows,babel}
\pgfplotsset{compat=1.18}
\usepgfplotslibrary{fillbetween}
\usepackage[np]{numprint}
\usepackage{listings,makecell,pas-tableur,twoopt}
\let\barre\undefined
\usepackage[pstricks]{tablvar}
\definecolor{darkgreen}{RGB}{0,100,0}
\providecommand{\red}{\color{red}}
\providecommand{\blue}{\color{blue}}
\providecommand{\green}{\color{green}}
\providecommand{\black}{\color{black}}
\providecommand{\white}{\color{white}}
\providecommand{\gray}{\color{gray}}
\providecommand{\cyan}{\color{cyan}}
\providecommand{\magenta}{\color{magenta}}
\providecommand{\yellow}{\color{yellow}}
\providecommand{\violet}{\color{violet}}
\providecommand{\orange}{\color{orange}}
\providecommand{\euro}{\texteuro}
\providecommand{\EUR}{\texteuro}
\providecommand{\ieme}{\textsuperscript{e}}
\providecommand{\degres}{\ensuremath{^\circ}}
\providecommand{\degree}{\ensuremath{^\circ}}
\providecommand{\vv}[1]{\overrightarrow{#1}}
\providecommand{\Coord}[1]{#1}
\providecommand{\decosix}{}
\providecommand{\starredbullet}{\ensuremath{\bullet}}
\pagestyle{empty}
\setlength{\parindent}{0pt}
\usepackage{adjustbox}\DeclareUnicodeCharacter{2236}{:}\DeclareUnicodeCharacter{2190}{\ensuremath{\leftarrow}}
\providecommand{\R}{\mathbb{R}}
\providecommand{\N}{\mathbb{N}}
\providecommand{\Z}{\mathbb{Z}}
\providecommand{\Q}{\mathbb{Q}}
\providecommand{\D}{\mathbb{D}}
\providecommand{\C}{\mathbb{C}}
\providecommand{\e}{{\mathrm{e}}}
\providecommand{\vect}[1]{\overrightarrow{#1}}
\providecommand{\vectt}[1]{\overrightarrow{\mathrm{#1}}}
\providecommand{\vvt}[1]{\overrightarrow{\mathrm{#1}}}
\providecommand{\np}[2][]{{#2~#1}}
\providecommand{\vv}[1]{{\overrightarrow{#1}}}
\renewcommand{\d}{{\mathrm{d}}}
\providecommand{\degres}{{^\circ}}
\providecommand{\degre}{{^\circ}}
\providecommand{\blue}{\color{blue}}
\providecommand{\red}{\color{red}}
\providecommand{\enskip}{\quad}
\providecommand{\llbracket}{\lbrack\!\lbrack}
\providecommand{\rrbracket}{\rbrack\!\rbrack}
\providecommand{\euro}{{\text{€}}}
\providecommand{\no}{{\text{n°}}}
\providecommand{\interval}[2]{{[#1;#2]}}
\providecommand{\intervFO}[2]{{[#1;#2[}}
\providecommand{\dotfill}{\ldots}

\providecommand{\ds}{\displaystyle}
\providecommand{\textding}[1]{#1}
\providecommand{\Oij}{\(\mathrm{O};\vec{\imath},\vec{\jmath}\)}
\providecommand{\Oijk}{\(\mathrm{O};\vec{\imath},\vec{\jmath},\vec{k}\)}
\providecommand{\Ouv}{\(\mathrm{O};\vec{u},\vec{v}\)}
\providecommand{\intervFF}[2]{\ensuremath{[#1;#2]}}
\renewcommand{\euro}{\ifmmode\text{\texteuro}\else\texteuro\fi}
\DeclareUnicodeCharacter{00B0}{\ifmmode^{\circ}\else\textdegree\fi}
\DeclareUnicodeCharacter{20AC}{\ifmmode\text{\texteuro}\else\texteuro\fi}
% Énoncés : Éducation nationale. Transcription et corrigés : APMEP, Denis Vergès et contributeurs. Découpage et classement : CoopMaths, Sébastien Lozano, Éric Elter et contributeurs. Adaptation de la présentation : Maths Explorer.
% CC BY-SA : https://coopmaths.fr/www/about/
% https://coopmaths.fr/alea/static/bac/2025/tex/bac_2025_06_sujet2_metropole_1.tex
\begin{document}
\noindent\textit{APMEP / CoopMaths. CC BY-SA. Présentation adaptée par Maths Explorer.}\par\medskip
\medskip

Les deux parties peuvent être traitées indépendamment.

Dans cet exercice on s'intéresse à des personnes venues séjourner dans un centre multisports au cours du week-end.

Les résultats des probabilités demandées seront arrondis au millième si nécessaire.

\bigskip

\begin{center}\textbf{Partie A} \end{center}

 Le centre propose aux personnes venues pour un week-end une fonmule d'initiation
au roller composée de deux séances de cours.

On choisit au hasard une personne parmi celles ayant souscrit à cette formule.

\medskip

On désigne par $A$ et $B$ les évènements suivants :

\begin{itemize}[label=$\bullet~$]
\item $A$ : \og La personne chute pendant la première séance\fg ;
\item $B$ : \og La personne chute pendant la deuxième séance \fg.
\end{itemize}

Pour un évènement $E$ quelconque, on note $P(E)$ sa probabilité et $\overline{E}$ son événement contraire.

Des observations permettent d'admettre que $P(A) = 0,6$.

De plus on constate que:

\begin{itemize}[label=$\bullet~$]
\item Si la personne chute pendant la première séance, la probabilité qu'elle chute pendant la deuxième est de $0,3$ ;
\item Si la personne ne chute pas pendant la première séance, la probabilité qu'elle chute pendant la deuxième est de $0,4$.
\end{itemize}

\medskip

\begin{enumerate}
\item Représenter la situation par un arbre pondéré.
\item Calculer la probabilité $P \left (\overline{A} \cap \overline{B} \right )$ et interpréter le résultat.
\item Montrer que $P(B) = 0,34$.
\item La personne ne chute pas pendant la deuxième séance de cours.

Calculer la probabilité qu'elle n'ait pas chuté lors de la première séance.
\item On appelle $X$ la variable aléatoire qui, à chaque échantillon de $100$ personnes ayant souscrit à la formule, associe le nombre d'entre elles n'ayant chuté ni lors de la première ni lors de la deuxième séance.

On assimile le choix d'un échantillon de 100 personnes à un tirage avec remise.

On admet que la probabilité qu'une personne ne chute ni lors de la première ni lors de la deuxième séance est de $0,24$.
	\begin{enumerate}
		\item Montrer que la variable aléatoire $X$ suit une loi binomiale dont on précisera les paramètres.
		\item Quelle est la probabilité d'avoir, dans un échantillon de $100$ personnes ayant souscrit à la formule, au moins $20$ personnes qui ne chutent ni lors de la première ni lors de la deuxième séance ?
		\item Calculer l'espérance $E(X)$ et interpréter le résultat dans le contexte de l'exercice.
	\end{enumerate}
\end{enumerate}

\begin{center}\textbf{Partie B} \end{center}

On choisit au hasard une personne venue  un week-end au centre multisport. On note $T_1$ la variable  aléatoire donnant son temps d'attente total en minute avant les accès aux activités sportives pendant la journée du samedi et $T_2$ la variable aléatoire donnant son temps d'attente total en minutes avant les accès aux activités sportives pendant la journée du dimanche.

On admet que :

\begin{itemize}[label=$\bullet~$]
\item $T_1$ suit une loi de probabilité d'espérance $E(T_1) = 40$ et d'écart-type $\sigma(T_1) = 10$ ;
\item $T_2$ suit une loi de probabilité d'espérance $E(T_2) = 60$ et d'écart-type $\sigma(T_2) = 16$ ;
\item  les variables aléatoires $T_1$ et $T_2$ sont indépendantes.
\end{itemize}

\medskip

On note $T$ la variable aléatoire donnant le temps total d'attente avant les accès au activités sportives lors des deux jours, exprimé en minute. Ainsi on a $T = T_1 + T_2$.

\medskip

\begin{enumerate}
\item Déterminer l'espérance $E(T)$ de la variable aléatoire $T$. Interpréter le résultat dans le contexte de l'exercice,
\item Montrer que la variance $V(T)$ de la variable aléatoire $T$ est égale à 356.
\item À l'aide de l'inégalité de Bienaymé-Tchebychev, montrer que, pour une personne
choisie au hasard parmi celles venues un week-end au centre multisports, la probabilité que son temps total d'attente $T$ soit strictement compris entre $60$ et $140$ minutes est supérieure à $0,77$.
\end{enumerate}

\bigskip

\end{document}
